\documentclass[11pt]{article}
\usepackage[margin=1in]{geometry}
\usepackage{amsmath,amssymb,booktabs,hyperref}
\usepackage{graphicx}
\usepackage{microtype}
\hypersetup{colorlinks=true,citecolor=blue,urlcolor=blue,linkcolor=blue}

\newcommand{\IIGBM}{\textsc{IIGBM}}
\newcommand{\E}{\mathbb{E}}

\title{Inherently Interpretable Gradient Boosting Machines:\\
Functional ANOVA, Model-Induced Kernels and Component Curvature}
\author{Nupur Krishnan \and Agus Sudjianto}
\date{\today}

\begin{document}
\maketitle

\begin{abstract}
Gradient boosted decision trees (GBDTs) predict well on tabular data but hide
how they combine features. We build a GBDT whose every explanatory object comes
from the fitted model itself. The workflow fits an additive GBDT with one
semantic feature group per tree, screens residual interactions with depth-two
trees and fits a deep residual model constrained to the screened pairs.
Functional ANOVA (fANOVA) reads off interaction order and component importance.
For each fANOVA component we build a positive semidefinite (PSD) kernel from the
component's own model responses and apply its graph Laplacian to the component
values, giving a local curvature that measures how sharply the learned effect
changes across inputs the model treats as similar. Importance says which learned
effects carry weight; curvature says whether an effect is broad and smooth or
localized and steep. We validate the curvature on planted interactions of known
shape and show it ranks a localized effect above a smooth one even when the
localized effect is eight times weaker. On California housing and bikeshare the
two quantities separate a dominant smooth spatial interaction from weak but
sharp calendar regimes. Every kernel and graph is induced by the fitted model;
no raw-feature similarity graph or surrogate is used.
\end{abstract}

\medskip
\noindent\textbf{Keywords:} inherently interpretable models; gradient boosting;
functional ANOVA; model-induced kernel; leaf co-membership kernel; graph-Laplacian
curvature; interaction detection; component importance.

\section{Introduction}

A fitted GBDT reports which features it splits on, but not how it combines them
into effects a domain expert can read. Global gain measures rank features
without describing the shape of an effect. Similarity graphs built from raw
input coordinates describe the data geometry, not the model's. We take a
model-centric route: the explanation, the kernel, the graph and the curvature
are all derived from the fitted GBDT.

The starting point is the fANOVA reading of tree ensembles. Tree depth bounds
the maximum interaction order, so a boosted ensemble can be read as a functional
ANOVA model rather than an opaque predictor \cite{Sudjianto2025,Hu2023}. A
depth-one ensemble is a generalized additive model (GAM); a depth-two ensemble
adds pairwise interactions; deeper trees express higher-order terms. Interaction
order is therefore a structural property of the model, available before any
attribution.

We add a geometric question about an identified fANOVA component: whether its
learned shape is broad and smooth or steep and localized. A PSD kernel built
from the component's own model responses, with a normalized graph-Laplacian
curvature of the component values, measures it. The measure is not causal and
is not Ollivier--Ricci curvature; it quantifies the local roughness of a known
fANOVA component. Section~\ref{sec:curvature} defines it and validates it against
interactions whose shape we control.

Contributions:
\begin{enumerate}
\item a hierarchical \IIGBM{} construction: an additive GBDT, a depth-two
residual interaction screen and a screened deep residual model;
\item PSD kernels induced by the deep model's fANOVA interaction-contrast feature
maps;
\item a component curvature from the component's kernel graph, validated on
planted smooth and localized interactions with a permutation null and stability
checks; and
\item studies on California housing and bikeshare that separate importance from
local shape.
\end{enumerate}

\section{Inherently Interpretable GBDT Construction}

\subsection{Stage I: additive tree ensemble}

Let $x=(x_1,\ldots,x_p)$ and let $y$ be the transformed response. We fit an
XGBoost ensemble in which each tree may split only within one semantic feature
group. A tree can use several splits of its group but cannot combine two groups,
so the ensemble is a GAM,
\begin{equation}
 f_1(x)=\beta_0+\sum_{j=1}^{p} f_j(x_j).
\end{equation}
Deep univariate trees fit flexible nonlinear main effects while additivity holds
exactly.

\subsection{Stage II: residual interaction screening}

To keep in-sample GAM artifacts from masquerading as interactions, we form
out-of-fold predictions $\hat f_1^{\mathrm{OOF}}$ and fit a depth-two GBDT to the
residual
\begin{equation}
 r^{\mathrm{OOF}}=y-\hat f_1^{\mathrm{OOF}}(x).
\end{equation}
In a depth-two tree the root and each child split give a candidate feature pair.
We aggregate child-split gains to rank pairs. This stage screens; it does not
supply the explanatory kernel.

\subsection{Stage III: screened deep residual model}

We refit the GAM on all training data and fit a deep residual model
\begin{equation}
 g_{\mathrm{deep}}(x)\approx y-f_1(x),
 \qquad
 f(x)=f_1(x)+g_{\mathrm{deep}}(x),
\end{equation}
with XGBoost interaction constraints drawn from the screened pairs. Overlapping
pairs are allowed, which lets the deep model compose higher-order terms where the
data support them. Stage II supplies an interpretable interaction prior; the
final deep model refines it.

\section{Functional ANOVA and Model-Induced Kernels}

\subsection{fANOVA components and importance}

Fix a reference measure $P_X$ with density $w$. Decompose the deep residual model
\begin{equation}
 g_{\mathrm{deep}}(x)=\sum_{S\subseteq\{1,\ldots,p\}}g_S(x_S),
\end{equation}
where $g_S$ is the fANOVA component on feature subset $S$. Components are centered
under the reference, $\E[g_S(X_S)]=0$, so component importance is the variance
\begin{equation}
 A_S=\mathrm{Var}\!\left(g_S(X_S)\right)=\E\!\left[g_S(X_S)^2\right].
\end{equation}
Centering matters in practice: without it a residual constant inflates
$A_S$. In our California study the Latitude$\times$Longitude contrast carried an
offset of roughly $0.9$ standard deviations; centering the component before
measuring $A_S$ cut its reported energy from $0.0136$ to $0.0075$ and left the
curvature unchanged (curvature is offset-invariant; Section~\ref{sec:curvature}).

We estimate components with a finite empirical background rather than an exact
purification. Given background draws $B^{(1)},\ldots,B^{(M)}$, define for a pair
$S=\{j,k\}$
\begin{equation}
 \Psi_{jk,m}(x)=g_{\mathrm{deep}}(x_j,x_k,B^{(m)}_{-jk})
 -g_{\mathrm{deep}}(x_j,B^{(m)}_{-j})
 -g_{\mathrm{deep}}(x_k,B^{(m)}_{-k})
 +g_{\mathrm{deep}}(B^{(m)}).
\label{eq:contrast}
\end{equation}
The component estimate is $\hat g_{jk}(x)=M^{-1}\sum_m\Psi_{jk,m}(x)$, centered
over the evaluation set. The reference measure is the joint empirical background,
so under correlated predictors the components are not exactly orthogonal; exact
tree-term aggregation and purification would remove the residual dependence
\cite{Hooker2007,Lengerich2020,Hu2023}.

\subsection{Hierarchical orthogonality}
\label{sec:ortho}

Fitting interactions on the GAM residual is what makes the importances behave
like a decomposition. The hierarchical-orthogonality proposition of Hu et
al.~\cite{Hu2023} states that if the lower-order terms minimize the weighted
projection $\int\!\big(f-\sum_j f_{-j}\big)^2 w\,dx$, the residual
$\tilde f=f-\sum_j f_{-j}$ is orthogonal to every lower-dimensional function
$h$: $\int \tilde f(x)\,h(x_{-j})\,w(x)\,dx=0$. The argument is variational: a
nonzero inner product would let a small multiple of $h$ lower the projection
loss, contradicting optimality. Stage III fits within this orthogonal
complement, so its interaction components are orthogonal to the main effects
under the reference measure by construction.

We check the finite-sample gap directly. On California housing, the empirical
correlation between each main effect and each pair component has mean $0.070$
across all pairs, and the pair components are mutually near-orthogonal
(mean $|{\rm corr}|=0.058$). One pair reaches $0.441$ against a main effect: the
correlated-feature regime where XGBoost does not solve the weighted projection
exactly. The mean near zero is consistent with the proposition; the residual
leakage marks where purification would help.

\subsection{From trees to kernels}
\label{sec:treekernel}

A fitted tree induces a kernel on its own, before any of the fANOVA machinery.
Route an input $x$ to its terminal leaf $\ell(x)$ and write the one-hot leaf
embedding $z(x)=e_{\ell(x)}\in\{0,1\}^L$, a vector with a single one in the
coordinate of the leaf $x$ lands in. The leaf co-membership kernel is the inner
product of these embeddings,
\begin{equation}
 K_T(x,x')=z(x)^\top z(x')=\mathbf 1\{\ell(x)=\ell(x')\},
\end{equation}
one when two points share a leaf and zero otherwise. Stacking $n$ points into a
matrix $Z$ with rows $z(x_i)^\top$ gives the Gram matrix $K=ZZ^\top$, so $K_T$ is
PSD by construction. An ensemble of $M$ trees stacks the per-tree embeddings into
a feature map
$\phi_{\mathrm{GBDT}}(x)=(\sqrt{\alpha_1}\,z_1(x)^\top,\ldots,\sqrt{\alpha_M}\,z_M(x)^\top)^\top$
and sums the co-membership kernels,
\begin{equation}
 K_{\mathrm{GBDT}}(x,x')=\sum_{t=1}^M \alpha_t\,z_t(x)^\top z_t(x'),
\end{equation}
which with $\alpha_t=1/M$ counts the fraction of trees that route two points to
the same leaf. Similarity is supervised and model-induced: two inputs are close
when the fitted ensemble repeatedly places them together, not when they are near
in raw coordinates \cite{SudjiantoZhang2026}.

\subsection{PSD component kernels}

The component kernel keeps this recipe --- a model-induced feature map, PSD as an
inner product --- but replaces the raw leaf indicators with the deep model's
interaction responses, so the geometry isolates one fANOVA component instead of
the full partition. The contrast vector
$\Psi_S(x)=(\Psi_{S,1}(x),\ldots,\Psi_{S,M}(x))^\top$ plays the role of the
feature map, and
\begin{equation}
 K_S(x,x')=\frac{1}{M}\Psi_S(x)^\top\Psi_S(x')
\label{eq:kernel}
\end{equation}
is PSD by the same Gram-matrix argument. Where the leaf-membership map $z(x)$
records which region $x$ occupies, $\Psi_S(x)$ records how the deep model's
prediction moves when features $S$ are set to $x$ against a background. Two points
are close under $K_S$ when that interaction response matches, not merely when they
share leaves.

\section{Component Curvature}
\label{sec:curvature}

\subsection{What the curvature measures, in words}

Each component $g_S$ gives every input one number: the size of the learned
interaction there. We ask whether that number changes gently or abruptly between
inputs the model treats as similar. Two ingredients make this precise. First,
``similar'' is the component's own kernel: keep, for each point, its $k$ nearest
neighbors in the feature-map space $\Psi_S(x)/\sqrt{M}$. Second, at each point
take the weighted average of its neighbors' component values and subtract it from
the point's own value. Where the effect is locally flat the point sits near its
neighborhood average and the difference is small; where the effect steps or
spikes (a threshold, a narrow peak), the point departs from its neighbors and
the difference is large. Squaring and averaging over points gives the roughness;
comparing it to the roughness of the same values shuffled at random over the
graph turns it into a scale that depends on neither the effect's amplitude nor
how many distinct states the component has. Readers who
know spectral graph theory will recognize the per-point quantity as the graph
Laplacian applied to the component, $c_S=(I-W)\,g_S$, and the summary as a
normalized Dirichlet energy \cite{Chung1997,BelkinNiyogi2003}.

\subsection{Definition}

Build a $k$-nearest-neighbor graph in $\Psi_S(x)/\sqrt{M}$ with normalized local
weights $w_{i\ell}$. The local curvature is
\begin{equation}
 c_S(x_i)=\hat g_S(x_i)-\sum_{\ell\in N(i)}w_{i\ell}\,\hat g_S(x_\ell),
\end{equation}
Write the roughness $R(g)=N^{-1}\sum_{i=1}^N c_S(x_i)^2$, weighting each graph
node by the number of points it collapses. The component summary normalizes it by
its mean under random permutations $\pi$ of the component values over the graph
nodes,
\begin{equation}
 Q_S=\frac{R(\hat g_S)}{\mathbb{E}_\pi\,R(\hat g_S\circ\pi)}.
\label{eq:curvature}
\end{equation}
The permutation null absorbs the graph's size and structure, so $Q_S$ compares a
component's roughness to what that same graph produces by chance and stays
comparable across components with very different numbers of distinct states. A
two-state binary feature sits on a two-node graph whose only permutation leaves
$R$ unchanged, so $Q_S=1$ rather than a tiny graph inflating it. $Q_S$ is
amplitude-invariant, since numerator and null scale together, and
centering-invariant, since a constant added to $g_S$ cancels in $c_S$ (the
neighbor weights sum to one). A value below $1$ is smoother than chance; a value
near $1$ carries no smoothness signal, either rough or on too few states to
resolve. Points whose feature maps coincide are collapsed to one node
before the neighbor search.

Importance alone cannot tell a broad effect from a narrow one. Two components
with equal $A_S$ can differ in whether their energy is spread across the input
space or concentrated at a few thresholds, and $A_S$ assigns both the same
number. Curvature resolves the ambiguity, so the pair $(A_S,Q_S)$ carries
information neither coordinate holds alone. High $A_S$ with low $Q_S$ is an
important broad effect, safely summarized by a simple trend. High $A_S$ with high
$Q_S$ is an important effect built from thresholds or regimes, worth plotting and
monitoring because a small input change can move the prediction sharply. Low
$A_S$ with high $Q_S$ is a weak effect on sparse or discrete support, to inspect
for data coverage or instability rather than act on. The reading is the same for
a main effect and an interaction; only the order of $S$ changes.

\subsection{Validation on planted interactions}
\label{sec:validation}

$Q_S$ is a descriptive statistic, so we test whether it ranks shape as intended
on data where the shape is known. We draw six independent uniform features and
build a response with two planted pair effects: a smooth bilinear term on
$(u_1,u_2)$ and a sharp Gaussian bump on $(u_3,u_4)$, plus additive main effects
and light noise. We run the full workflow and measure $A_S$ and $Q_S$ on held-out
points.

Table~\ref{tab:validation} reports the result. The localized bump has eight times
less variance than the smooth term ($A=0.005$ versus $0.043$) yet more than twice
the curvature ($Q=0.0040$ versus $0.0015$), the low-importance, high-curvature
pattern the method is meant to surface. Because $Q_S$ already divides by the
permutation null, both values sit far below $1$, and the permutation $z$-scores
($-79$ and $-19$) confirm each component is much smoother than a random labeling
rather than noise. The ranking is stable across five data seeds
($Q=0.002\pm0.000$ smooth, $0.005\pm0.000$ localized) and preserved across
$k\in\{10,15,25\}$.

\begin{table}[ht]
\centering
\caption{Curvature validation on planted interactions (held-out points).}
\label{tab:validation}
\begin{tabular}{lrrr}
\toprule
Planted effect & Importance $A_S$ & Curvature $Q_S$ & Null $z$\\
\midrule
Smooth bilinear $(u_1,u_2)$ & 0.0433 & 0.0015 & $-79$\\
Localized bump $(u_3,u_4)$  & 0.0052 & 0.0040 & $-19$\\
\bottomrule
\end{tabular}
\end{table}

Two limits follow from the same experiment. The dynamic range is modest: a sharp
continuous bump moves $Q_S$ by a factor of about two, not orders of magnitude,
because a fANOVA component is already a smooth functional over its own feature-map
graph. And $Q_S$ reads on one scale: a value near the chance ceiling $Q_S\approx1$
means the component is either rough or carried by too few states to
resolve, so compare components rather than reading an absolute threshold.

\section{Experiments}

\subsection{Protocol}

We study California housing and UCI hourly bikeshare. Both positive responses are
transformed by $y=\log(1+\text{target})$. Screening uses five-fold out-of-fold
GAM residuals. The final GAM, the deep residual model and all component kernels
are fit or evaluated without held-out labels. The finite-background fANOVA
estimate uses $96$ reproducible training-background draws. Results reproduce under
XGBoost 3.0 with pandas 2.2; later XGBoost with pandas 3.0 drops DataFrame
feature names, so interaction constraints must be passed as integer indices.

\subsection{California housing}

The main-effect GAM reaches a held-out log-scale MAE of $0.1264$. The screened
deep residual model improves it to $0.1028$; an unrestricted depth-six XGBoost
reaches $0.0934$. The staged model recovers most of the nonadditive structure
while keeping an explicit main-effect and interaction hierarchy.

Table~\ref{tab:housing} reports deep-model pair components with centered $A_S$.
Latitude$\times$Longitude dominates the ANOVA energy but has the lowest
curvature, consistent with a broad smooth spatial surface. The MedInc pairs are
far weaker but more localized: MedInc$\times$Population carries a fraction of the
spatial energy yet an order of magnitude more curvature.

\begin{table}[ht]
\centering
\caption{California housing: deep-model fANOVA components (centered $A_S$).}
\label{tab:housing}
\begin{tabular}{lrr}
\toprule
Interaction & Importance $A_S$ & Curvature $Q_S$\\
\midrule
Latitude $\times$ Longitude & 0.00747 & 0.01434\\
MedInc $\times$ AveRooms    & 0.00092 & 0.03689\\
MedInc $\times$ AveOccup    & 0.00032 & 0.06048\\
MedInc $\times$ Population   & 0.00030 & 0.15196\\
MedInc $\times$ HouseAge    & 0.00021 & 0.09240\\
\bottomrule
\end{tabular}
\end{table}

\subsection{Bikeshare}

For hourly bikeshare the response is $\log(1+\texttt{cnt})$. The GAM reaches a
held-out log-scale MAE of $0.4409$; the screened deep residual model reaches
$0.2761$. Table~\ref{tab:bike} separates importance from shape.
Hour$\times$weekday is the most important interaction; hour$\times$workingday is
smaller but twice as curved, matching the abrupt commute-hour switch between
working and non-working days. Month$\times$year is weak in energy yet the most
curved pair ($Q=0.71$), a sparse seasonal adjustment that varies year to year.
The support-invariant curvature also corrects a reading the raw statistic got
wrong: holiday$\times$month looked highly curved under a variance normalization,
but against its own permutation null it sits at $Q=0.055$, so it is weak and
smooth rather than a sharp regime. Its earlier high value came from a 21-state
graph, the small-support inflation the null normalization removes.

\begin{table}[ht]
\centering
\caption{Bikeshare: deep-model fANOVA components (centered $A_S$).}
\label{tab:bike}
\begin{tabular}{lrr}
\toprule
Interaction & Importance $A_S$ & Curvature $Q_S$\\
\midrule
hour $\times$ weekday    & 0.03862 & 0.19724\\
hour $\times$ workingday & 0.02392 & 0.39681\\
humidity $\times$ windspeed & 0.00560 & 0.05957\\
hour $\times$ temperature & 0.00468 & 0.05311\\
month $\times$ year      & 0.00189 & 0.70816\\
holiday $\times$ month   & 0.00086 & 0.05493\\
\bottomrule
\end{tabular}
\end{table}

\subsection{Main effects}

Importance and curvature apply to every fANOVA component, not only interactions.
Tables~\ref{tab:housing-main} and \ref{tab:bike-main} report the first-order
components of the full model $f=f_1+g_{\mathrm{deep}}$. Two readings carry over
from the interactions, and one caveat becomes concrete.

Main effects hold most of the explained variation. On California the three
leading main effects (MedInc, Latitude, Longitude) have importance
$0.02$--$0.04$, an order of magnitude above the strongest interaction
(Latitude$\times$Longitude at $0.0075$); on bikeshare the hour-of-day effect
alone has importance $1.30$, far above any interaction. The additive backbone
dominates, as the staged construction intends.

Among the leading effects, low curvature marks broad dependence and higher
curvature marks structure. California's MedInc, Latitude and Longitude are smooth
($Q\le 0.011$), consistent with monotone price gradients; bikeshare's hour-of-day
effect is the largest and also the most curved of its continuous-scale features
($Q=0.39$), matching the twin commute peaks a smooth trend would miss, while
temperature is smooth ($Q=0.011$). The weak effects are the curved ones:
California's HouseAge carries almost no energy ($A=0.0001$) yet $Q=0.17$, a sparse
local wiggle to inspect for data support rather than a broad effect to act on.

The two-state calendar features (year, holiday, workingday) show the support
normalization at work. Each sits on a two-node graph whose single permutation
leaves the roughness unchanged, so $Q_S$ pins to the chance value $1$ instead of
the inflated $Q>30$ a variance normalization gave holiday. A value at the ceiling
reflects that two states cannot express curvature, not a spurious spike.

\begin{table}[ht]
\centering
\caption{California housing: main-effect importance and curvature (centered).}
\label{tab:housing-main}
\begin{tabular}{lrrr}
\toprule
Feature & Importance $A_j$ & Curvature $Q_j$ & States\\
\midrule
Latitude   & 0.04398 & 0.01145 & 208\\
Longitude  & 0.04348 & 0.00894 & 236\\
MedInc     & 0.02327 & 0.00621 & 256\\
AveOccup   & 0.00507 & 0.00438 & 256\\
AveRooms   & 0.00162 & 0.03043 & 255\\
AveBedrms  & 0.00025 & 0.05422 & 252\\
HouseAge   & 0.00012 & 0.17430 & 52\\
Population & 0.00008 & 0.07127 & 212\\
\bottomrule
\end{tabular}
\end{table}

\begin{table}[ht]
\centering
\caption{Bikeshare: main-effect importance and curvature (centered). Two-state
features (year, holiday, workingday) sit on degenerate two-node graphs and pin to
the chance value $Q_j=1$.}
\label{tab:bike-main}
\begin{tabular}{lrrr}
\toprule
Feature & Importance $A_j$ & Curvature $Q_j$ & States\\
\midrule
hour        & 1.30207 & 0.38958 & 24\\
temperature & 0.09661 & 0.01071 & 48\\
year        & 0.04690 & 1.00000 & 2\\
month       & 0.01766 & 0.83008 & 12\\
weathersit  & 0.01564 & 0.71586 & 4\\
humidity    & 0.01314 & 0.10098 & 80\\
windspeed   & 0.00500 & 0.11322 & 26\\
weekday     & 0.00399 & 0.93946 & 7\\
holiday     & 0.00111 & 1.00000 & 2\\
workingday  & 0.00020 & 1.00000 & 2\\
\bottomrule
\end{tabular}
\end{table}

\section{Discussion}
\label{sec:disc}

The method keeps three questions apart. Tree depth and fANOVA determine
interaction order. Centered ANOVA energy measures global importance. Curvature
describes local shape, and it augments importance rather than standing in for it.

One caution remains from the construction. The permutation null makes $Q_S$
comparable across components of different support cardinality, and a two-node
binary graph pins to $Q_S=1$ rather than inflating, but the null cannot
manufacture resolution the support does not contain. A component on very few
states reports the chance value $1$, which says curvature is unresolved there, not
that the effect is rough. Read a $Q_S$ near $1$ as no signal, and rely on it as a
shape measure only for components with enough distinct states to populate the
graph.

The method explains the fitted model, not the data-generating mechanism, and
should not be read causally. The reference measure sets what ``importance'' and
``orthogonality'' mean; we use the joint empirical background and state it,
because under correlated predictors a product-measure decomposition would differ.

\section{Conclusion}

\IIGBM{} runs from tree structure to interaction order, PSD component kernels, and
local curvature, all induced by the fitted model. fANOVA identifies which
interaction the model uses; component curvature, validated against interactions
of known shape, describes how that interaction behaves locally.

\begin{thebibliography}{9}
\bibitem{Sudjianto2025}
A. Sudjianto.
\newblock From Classical ANOVA to Gradient Boosted Trees: A Functional ANOVA
Perspective on Inherently Interpretable Models.
\newblock SSRN Working Paper 5315091, 2025.
\newblock \url{https://papers.ssrn.com/sol3/papers.cfm?abstract_id=5315091}.

\bibitem{Hu2023}
L. Hu, V. N. Nair, A. Sudjianto, A. Zhang, and J. Chen.
\newblock Interpretable Machine Learning based on Functional ANOVA Framework:
Algorithms and Comparisons.
\newblock arXiv:2305.15670, 2023.

\bibitem{SudjiantoZhang2026}
A. Sudjianto and A. Zhang.
\newblock Generalized Nadaraya--Watson Operators for Learned Kernel Geometry: A
Unified View of Gradient-Boosted Trees, Kernel Regression and Self-Attention.
\newblock SSRN Working Paper 6823639, 2026.
\newblock \url{https://papers.ssrn.com/sol3/papers.cfm?abstract_id=6823639}.

\bibitem{Chung1997}
F. R. K. Chung.
\newblock \emph{Spectral Graph Theory}.
\newblock American Mathematical Society, 1997.

\bibitem{BelkinNiyogi2003}
M. Belkin and P. Niyogi.
\newblock Laplacian Eigenmaps for Dimensionality Reduction and Data
Representation.
\newblock \emph{Neural Computation}, 15(6):1373--1396, 2003.

\bibitem{Hooker2007}
G. Hooker.
\newblock Generalized Functional ANOVA Diagnostics for High-Dimensional
Functions of Dependent Variables.
\newblock \emph{Journal of Computational and Graphical Statistics},
16(3):709--732, 2007.

\bibitem{Lengerich2020}
B. Lengerich, S. Tan, C.-H. Chang, G. Hooker, and R. Caruana.
\newblock Purifying Interaction Effects with the Functional ANOVA: An Efficient
Algorithm for Recovering Identifiable Additive Models.
\newblock In \emph{AISTATS}, 2020.

\bibitem{Friedman2001}
J. H. Friedman.
\newblock Greedy Function Approximation: A Gradient Boosting Machine.
\newblock \emph{The Annals of Statistics}, 29(5):1189--1232, 2001.

\bibitem{ChenGuestrin2016}
T. Chen and C. Guestrin.
\newblock XGBoost: A Scalable Tree Boosting System.
\newblock In \emph{Proceedings of KDD}, 2016.
\end{thebibliography}

\end{document}
